\documentclass[12pt]{article}

\usepackage[a4paper, total={6in, 8in}]{geometry}
\usepackage{textcomp}

\begin{document}
\setlength{\parindent}{0pt}

\def \t {\textrightarrow}
\def \v {\vspace{1em}}
\def \bi {\begin{itemize}}
\def \ei {\end{itemize}}
\def \s[#1] {\section*{#1}}
\def \ss[#1] {\subsection*{#1}}
\def \sss[#1] {\subsubsection*{#1}}

\s[Italo Svevo]
\ss[La coscienza di zeno]
I personaggi credono di avere in pugno la realta, ma vengono subissati in quanto si mostra una critica forte nei confronti di zeno cosini, che si dimostra inconcludente \\
Zeno cosini rinuncia alla terapia con surrentizie motivazioni \t il romanzo è sfaldato, l'impianto manzoniano è distorto \t il narratore è il protagonista, la focalizzazione premia il monologo interiore con discorso diretto e indiretto e indiretto libero \\
Il tempo dell'aziione non è lineare, ma segue i flussi della coscienza \t attenuati rispetto a Joyce, e l'ambiente presentato è filtrato dagli occhi di zeno cosini \\
Noi conosciamo attraverso i suoi pensieri e i suoi giudizi, attraverso i suoi occhi \t l'ambiente è quanto basta per traslarlo sulla pagina \\
I personaggi, nel caso dell'opera, sono correlati a una rete familiare:
\bi
  \item sorelle
  \item cognato guido
  \item e poi ci sono dottor S., che sta per sigmund freud
\ei
Gia Boccaccio \t elegia a madonna fiammetta è il primo romanzo psicologico \t e cornice rimanda a boccacio, e il suo decamerone (storie introdotte da una cornice) \\
Qua la cornice è una dichiarazione di intenti \\ 
È zeno che rifiuta la terapia \t qui c'è la morte dell'anima \t l'analista esce con un volto degradato della sua personalita \\
Dimostra la fallacia della psicanalisi \t zeno cosini prende una decisione di fatto, nella sua indecisione \t ha superato se stesso, ha deciso di interrompere la terapia \\
Il dottor S lo implora invece e ne fa una questione di denaro e di notorieta \t questo viene vissuto con il coniato ?? 

\v

Entrano le parole memoriale, autoinganno \t e soprattutto alibi \t la psicanalisi esce svilita \\
La letteratura mostra la svolta della sua decisione, ma la psicanalisi viene descritta da svevo come una disciplina che serve di piu ai poeti che ai pazienti \\
La parola chiave è la parola alibi = escamotage \\
Si trovano le singole porzioni di elaborazione di zeno, che sfalda la regolarita della sintassi \t è il linguaggio catastrofico della delusione, e deluso e disilluso per quanto riguarda i rapporti familiari \\
Sensi di colpa, autoconsapevolezza sempre sfalsata \t nel vizio del fuomo zeno da colpa al padre, si deresponsabilizza \t incolpare un terzo non deve esimerci dalle nostre colpe \\
Lo spostare la responsabilita è tipico dell inettitudine, che di fatto è una malattia (dell'anima) \\
Zeno e Augusta, sorella Malpeliti???? che è il ritratto della sanità \t mentre zeno è il diverso, l inetto \t ma alla fine lei è malata \\
Il cognato è invece ciò che rappresenta cio che lui non potra mai essere \t questo porta a perdersi, e alla frustrazione e delusione \\
Tutto parte dall'ottica del soggetto che guarda il mondo, siamo nel 24 \t la teoria di einstein è in quegli anni \t e anche svevo viene influenzato dalla diffusione di queste idee

\s[Umberto Saba]
Anche qua trieste, terra di confine, influenza asburgica, poliglottismo \\
Si chiama in realtà Umberto Coli, in arte Saba \t la critica di Grosser e Cataldi come "psicanalitico prima della psicanalisi" \\
È il sentire di un artista \\
Se svevo parte da un autobiografismo parziale (ha amore per letteratura ma non psicanalisi) \t sabba invece vive il dramma che non parla di psicanalisi in modo diretto \\
Il metauro di tasso ritorna, i lutti ingiustificati, le morti non conosciute (lui pensava la madre fosse morta ma era a roma) \\
Nella sua semplicita, metrica quasi inesistente, continuatore delle dinamiche petrarchesche del canzoniere, permette di avvicinarsi a una cultura europea e alla componente austrungarica della cultura del tempo \\
Legge freud prima che si diffondesse in italia \t l opera di freud non è accessibile negli anni venti, e qua ci sotto pone a una terapia psicanalitica con lo specialista vicino a freud \\
Nella brevitas i suoi testi hanno una pensosità che sfiora il mondo degli aforismi = piccole riflessioni pregne di significato \\
Si coglie il canto dell'io, l'inconoscio, nel romano "Ernesto" si parla di un adolescente omosessuale \\
Ha un amore particolare per la figlia, e fa anche lirica su di lei

\ss[Vita]
Nel 1883 nasce, da una nobile famiglia venenziana di origine ebraica \t famiglia di commercianti e borghesi \\
La famiglia con la madre che gli viene tolta, per cui viene affidato a una balia slovena: Peppa Sabaz \\
La balia è importante \t lo allatta, e questo instaura un legame fortissimo \\
Ha educazione repressiva da parte del padre, e un trauma doloroso che lo segna (la madre vera si allontana ed è malata psicologicamente) lo porta a vivere un dualismo \\
Lui incolpa il padre, che definisce l'assassino \\
Incappa in una nevrosi che lo accompagna fino alla morte \t sposa pero Lina (Carolina Befler), per la quale compone molto \t poi hanno bambina di nome Linuccia \\
Lui dedica tutto se stesso alla figlia \\
Lo tocca la politica, e deve andare al fronte (intorno al 1915) \t poi torna a trieste, dove compra libreria antiquaria \t ed è la passione di librario e di scrittura che trovano compimento nella sua opera "Canzoniere" \\
Il canzoniere di saba raccogliere le sue liriche in ordine diaristico, secondo un impianto per cui l autore parte dalla sua esperienza e ne parla \\
Dal 29 inizia cura con allievo di Freud, Weiss \t irrisolti, verbalizzazione \t fino al 31 \\
Nel 31 inizia a rifletter sull inefficacia della terapia \t non si risolve la nevrosi \\
Lui si fa autoanalisi \t la riflessione è furtto di uno scavo dentro di se \\
In questi anni ci si avvicina alle leggi fascistissime \t e le persecuzioni raziali \\
Si rifugia a firenze, dove incontra Montale e al bar dove si trovano, discutano delle questioni attuali \\
Alterna momenti di depressione e momenti di serenita \\
Nel dopo guerra la sua poesia viene riconosciuta, come valida linea alternativa a quella novecentista \\
Il legame con la civilta austriaca, ed è vorace del conoscere \t studia anche Freud e Nietzsche \t la riflessione antiaccademica di Freud e l'inconsistenza della morale nietzschiana \\
La trama, e l'assenza di una morale \t che viene letta attraverso il principio primo del piacere \t che fa lecito cio che non lo era percepito \\
L'autoedonismo \t è affermazione di se in Nietzsche

\v

War poets \t i poeti della guerra, chi sono ??? \\
Il poeta per lui si dedicare allo scandaglio interiore (seneca e petrarca anche) \t se devo scandagliare la mia anima, la devo percepire come un abisso \t non siamo lontani dal poeta di Ungaretti, che è il palombaro \\
Il poeta è palombaro che fa emergere la verità dai fondali \\
La sua poesia è chiamata la poetica dell'onestà \t il poeta non è diverso dagli altri uomini \t ognuno ha i propri segreti che non si possono tirare fuori \\
Il nostro animo non si svela in toto, neanche a noi stessi e non solo agli altri \t il poeta e la poesia sono quello che gli ermetici ci sveleranno, ovvero qualcosa riposto nel profondo dell'animo \\
La poesia ha valore liberatorio, ed è un impegno con noi stessi che è dentro arrovellato nella psiche umana \t gli uomini di tutti i giorni non sono consapevoli di questo \\
Irrisolto, che grazie alla poesia diventa visibile

\v

Il canzoniere ha +400 testi, composto in mezzo secolo \\
Prima edizione nel 1931?, mentre la seconda nel 1945 (einaudi) \t edizione postuma nel 61 \t gestazione che accompagna mezzo secolo \\
Le poesie sono divise in tre volumi, a loro volta divise in 25 sezioni \t impianto rigido, ma le opere godono di una fluidita \\
Il confronto con petrarca presenta un canzoniere con testi con struttura organica e ordinata, mentre qua i testi seguono fatti, a volte collocati in modo arbitrario \\
Ci sono i versi dell'adolescenza \t cose di campagna \t trieste è una donna \t serena disperazione (ossimoro) \t cose leggere e vaganti \t l'amorosa spina (inclinazione ossimoro, o meglio alla psicomachia = dualismo interiore, fluctatio animae)
\end{document}
